\documentclass[11pt]{article}
\usepackage[margin=1in]{geometry}
\usepackage{xcolor}
\usepackage{enumitem}
\usepackage{hyperref}
\usepackage{booktabs}
\usepackage{amsmath}

\newcommand{\point}[1]{\paragraph{#1}}
\newcommand{\response}{\textcolor{blue}{\textbf{Response:}}~}
\newcommand{\lines}[1]{\texttt{[L#1]}}

\title{Author Response to Reviewer u34v\\(Rating: 4, Borderline Accept)}
\author{Anonymous Authors}
\date{}

\begin{document}
\maketitle

We thank Reviewer u34v for the detailed and constructive suggestions on presentation and clarity. Below we address each point in detail. All changes in the revised manuscript are marked in \textcolor{blue}{blue text}; line references (e.g., \lines{42--44}) refer to the revised \texttt{main.tex}.

\point{R3.1: Manuscript feels too fragmented due to titled paragraphs.}

\response We have revised the Discussion (Section~6, \lines{594--623}) to reduce fragmentation. Bold titled paragraphs (``An empirical property taxonomy,'' ``Relationship to prior work,'' ``Practical implications'') have been converted to flowing prose with topic sentences that integrate the content more naturally. Specifically: the taxonomy opening at \lines{605} now uses a topic sentence rather than a bold header; the prior work transition at \lines{621} flows naturally from the preceding text; and the practical implications at \lines{623} open with a transitional phrase. The section now reads as a cohesive narrative rather than a series of independent blocks.

\textbf{Specific changes:}
\begin{itemize}[itemsep=2pt]
    \item Taxonomy opening rewritten as flowing prose: \lines{605}
    \item Property categories (solvation, steric, electronic): \lines{607--611}
    \item Performance hierarchy rewritten with flowing transition: \lines{613}
    \item Prior work transition to flowing prose: \lines{621}
    \item Practical implications transition: \lines{623}
\end{itemize}

\point{R3.2: Research questions formulated as actual questions feels too informal.}

\response We have rephrased RQ1--RQ3 in the Introduction (\lines{147--149}) as declarative objectives rather than interrogative questions:
\begin{itemize}[itemsep=2pt]
    \item RQ1 \lines{147}: ``\textbf{Property taxonomy}: We characterize which molecular property types benefit from 3D conformer geometry and provide a mechanistic basis for the observed selectivity.''
    \item RQ2 \lines{148}: ``\textbf{Representation hierarchy}: We establish how pre-computed conformer ensemble methods compare to 2D fingerprints and end-to-end 3D GNNs (SchNet, PaiNN) across task types.''
    \item RQ3 \lines{149}: ``\textbf{Representation complexity}: We test whether complex covariance representations outperform simple summary statistics, grounded in a bias-variance analysis of parameterization complexity versus dataset size.''
\end{itemize}

\textbf{Specific changes:}
\begin{itemize}[itemsep=2pt]
    \item RQ1 rephrased as declarative objective: \lines{147}
    \item RQ2 rephrased as declarative objective: \lines{148}
    \item RQ3 rephrased as declarative objective: \lines{149}
\end{itemize}

\point{R3.3: Performance hierarchy is confused and could benefit from clearer presentation.}

\response We have substantially rewritten the performance hierarchy presentation in Section~6 at \lines{613}. The revised text walks through each of the four tiers with explicit evidence and, crucially, explains the progression from tier 4 to tier 1 in terms of how much 3D relational structure is preserved: ``DKO discards it into summary statistics, engineered descriptors preserve selected physical invariants, and GNNs retain the full atomic graph.'' This mechanistic framing---absent from the original---makes the logic of the hierarchy explicit rather than presenting it as an empirical ranking without explanation. The key result---that FP+3D descriptors match SchNet on ESOL (1.000 vs.\ 1.004)---is highlighted as evidence that carefully engineered features can approach end-to-end learning.

\textbf{Specific changes:}
\begin{itemize}[itemsep=2pt]
    \item Four-tier hierarchy rewritten with mechanistic explanation: \lines{613}
    \item Tier progression explained (summary statistics $\to$ physical invariants $\to$ full atomic graph): \lines{613}
    \item Key FP+3D vs.\ SchNet comparison highlighted: \lines{613}
    \item Abstract hierarchy description updated: \lines{82}
    \item Contribution 3 hierarchy description updated: \lines{157}
    \item Conclusion hierarchy summary: \lines{630}
\end{itemize}

\point{R3.4: Some text too direct; expand into more technical discourse.}

\response We have expanded terse statements throughout the manuscript into more technical language. Key locations:
\begin{itemize}[itemsep=2pt]
    \item Abstract performance hierarchy: \lines{82}---now specifies ``continuous-filter and equivariant message passing'' and ``pre-computed feature bottleneck''
    \item Discussion bias-variance analysis: \lines{617}---explicit parameter-count calculation (${\sim}180$ vs.\ ${\sim}7$ samples per parameter), connection to model selection theory
    \item Hierarchy explanation: \lines{613}---mechanistic framing of tier progression with evidence citations
    \item Practical decision framework: \lines{623}---expanded with specific thresholds and conditions
    \item Conclusion: \lines{625--633}---tightened with precise statistical claims and explicit practical recommendations
    \item GNN baseline justification: \lines{460}---technical comparison of continuous-filter vs.\ equivariant paradigms
    \item Truncation concern: \lines{515}---technical argument about curse of dimensionality vs.\ information loss
\end{itemize}

\vspace{1em}
\noindent\rule{\textwidth}{0.4pt}
\vspace{0.5em}

\noindent\textbf{Summary of changes relevant to Reviewer u34v:}
\begin{itemize}[itemsep=2pt]
    \item Rewrote Discussion for clarity and flow \lines{605}, \lines{613}, \lines{621}, \lines{623}
    \item Rephrased research questions as declarative objectives \lines{147--149}
    \item Expanded performance hierarchy with mechanistic explanation \lines{613}
    \item Expanded technical discourse throughout \lines{82}, \lines{460}, \lines{515}, \lines{617}
    \item All changes marked in \textcolor{blue}{blue text}
\end{itemize}

\end{document}
